
\documentclass{article}
\usepackage{colortbl}
\usepackage{makecell}
\usepackage{multirow}
\usepackage{supertabular}

\begin{document}

\newcounter{utterance}

\centering \large Interaction Transcript for game `codenames', experiment `risk\_low', episode 0 with Qwen3.6{-}35B{-}A3B{-}FP8{-}without{-}reasoning.
\vspace{24pt}

{ \footnotesize  \setcounter{utterance}{1}
\setlength{\tabcolsep}{0pt}
\begin{supertabular}{c@{$\;$}|p{.15\linewidth}@{}p{.15\linewidth}p{.15\linewidth}p{.15\linewidth}p{.15\linewidth}p{.15\linewidth}}
    \# & $\;$A & \multicolumn{4}{c}{Game Master} & $\;\:$B\\
    \hline

    \theutterance \stepcounter{utterance}  
    & & & \multicolumn{2}{p{0.3\linewidth}}{
        \cellcolor[rgb]{0.9,0.9,0.9}{
            \makecell[{{p{\linewidth}}}]{
                \texttt{\tiny{[GM$|$GM]}}
                {'hidden': {'team': ['революція', 'дошка', 'курча', 'відьма', 'сила', 'міст', 'клен', 'контракт', 'ліжко'], 'innocent': ['лід', 'лепрекон', 'екран', 'лікарня', 'бар', 'москва', 'вентилятор', 'машина'], 'opponent': ['злодій', 'лимон', 'залізо', 'прибулець', 'концерт', 'басейн', 'африка', 'стопа'], 'assassin': []}, 'revealed': {'team': {'team': [], 'innocent': [], 'opponent': [], 'assassin': []}, 'opponent': {'team': [], 'innocent': [], 'opponent': [], 'assassin': []}}}
            }
        }
    }
    & & \\ \\

    \theutterance \stepcounter{utterance}  
    & & \multicolumn{4}{p{0.6\linewidth}}{
        \cellcolor[rgb]{0.9,0.9,0.9}{
            \makecell[{{p{\linewidth}}}]{
                \texttt{\tiny{[P1$\langle$GM]}}
                \texttt{Надайте одну однословну підказку, яка пов'язана з деякими словами вашої команди зі списку нижче. Ви можете обрати, зі скількома і з якими саме словами пов'язана ваша підказка.} \\
\texttt{Мета полягає в тому, щоб ваш партнер відгадав усі слова вашої команди на основі підказок, які ви йому даєте, раніше, ніж інша команда відкриє всі свої. Якщо ваш партнер відгадає слово «суперника» або «відволікача», його хід відгадування закінчується, і команда суперників отримує право зробити хід. Якщо ваш партнер відгадає слово «вбивці», гра закінчується, і ваша команда негайно програє.} \\
\texttt{Чим із більшою кількістю слів пов'язана ваша підказка, тим швидше ви можете виграти гру, але тим складнішими можуть бути асоціації для відгадування вашим партнером, тож обирайте слово{-}підказку мудро. Слово{-}підказка має бути семантично пов'язаним із цільовими словами; воно не може бути одним зі слів у списках або містити їхні частини.} \\
\\ 
\texttt{Завжди надавайте вашу однословну підказку та ваш розділений комами список пов'язаних цільових слів у такому форматі й робіть відповіді якомога коротшими; ніколи не додавайте жодного іншого тексту, окрім того, що потрібно в цій формі:} \\
\\ 
\texttt{ПІДКАЗКА: <СЛОВО>} \\
\texttt{ЦІЛІ: <ЦІЛІ>} \\
\\ 
\texttt{Слова вашої команди: революція, дошка, курча, відьма, сила, міст, клен, контракт, ліжко.} \\
\texttt{Слова команди суперника: злодій, лимон, залізо, прибулець, концерт, басейн, африка, стопа.} \\
\texttt{Слова{-}відволікачі: лід, лепрекон, екран, лікарня, бар, москва, вентилятор, машина.} \\
\texttt{Слова{-}вбивці: .} \\
            }
        }
    }
    & \\ \\

    \theutterance \stepcounter{utterance}  
    & \multicolumn{4}{p{0.6\linewidth}}{
        \cellcolor[rgb]{0.9,0.9,0.9}{
            \makecell[{{p{\linewidth}}}]{
                \texttt{\tiny{[P1$\rangle$GM]}}
                \texttt{ПІДКАЗКА: будівля} \\
\texttt{ЦІЛІ: дошка, ліжко, міст} \\
            }
        }
    }
    & & \\ \\

    \theutterance \stepcounter{utterance}  
    & & & \multicolumn{2}{p{0.3\linewidth}}{
        \cellcolor[rgb]{0.9,0.9,0.9}{
            \makecell[{{p{\linewidth}}}]{
                \texttt{\tiny{[GM$|$GM]}}
                \texttt{будівля} \\
            }
        }
    }
    & & \\ \\

    \theutterance \stepcounter{utterance}  
    & & & \multicolumn{2}{p{0.3\linewidth}}{
        \cellcolor[rgb]{0.9,0.9,0.9}{
            \makecell[{{p{\linewidth}}}]{
                \texttt{\tiny{[GM$|$GM]}}
                ['дошка', 'ліжко', 'міст']
            }
        }
    }
    & & \\ \\

    \theutterance \stepcounter{utterance}  
    & & & \multicolumn{2}{p{0.3\linewidth}}{
        \cellcolor[rgb]{0.9,0.9,0.9}{
            \makecell[{{p{\linewidth}}}]{
                \texttt{\tiny{[GM$|$GM]}}
                {'word': 'дошка', 'assignment': 'team'}
            }
        }
    }
    & & \\ \\

    \theutterance \stepcounter{utterance}  
    & & & \multicolumn{2}{p{0.3\linewidth}}{
        \cellcolor[rgb]{0.9,0.9,0.9}{
            \makecell[{{p{\linewidth}}}]{
                \texttt{\tiny{[GM$|$GM]}}
                {'word': 'ліжко', 'assignment': 'team'}
            }
        }
    }
    & & \\ \\

    \theutterance \stepcounter{utterance}  
    & & & \multicolumn{2}{p{0.3\linewidth}}{
        \cellcolor[rgb]{0.9,0.9,0.9}{
            \makecell[{{p{\linewidth}}}]{
                \texttt{\tiny{[GM$|$GM]}}
                {'word': 'міст', 'assignment': 'team'}
            }
        }
    }
    & & \\ \\

    \theutterance \stepcounter{utterance}  
    & & \multicolumn{4}{p{0.6\linewidth}}{
        \cellcolor[rgb]{0.9,0.9,0.9}{
            \makecell[{{p{\linewidth}}}]{
                \texttt{\tiny{[GM$\rangle$P2]}}
                \texttt{Надайте розділений комами список із не більш ніж 3 слів із наведеного нижче списку, які найкраще пов'язані або найбільш тісно асоціюються зі словом «будівля». Завжди починайте свій список здогадок із ВІДГАДКА:  і не додавайте жодного іншого тексту у своїй відповіді.} \\
\\ 
\texttt{ліжко, сила, міст, відьма, клен, стопа, машина, бар, лепрекон, басейн, революція, лікарня, залізо, концерт, курча, лід, злодій, вентилятор, контракт, африка, москва, екран, лимон, дошка, прибулець} \\
            }
        }
    }
    & \\ \\

    \theutterance \stepcounter{utterance}  
    & & & \multicolumn{4}{p{0.6\linewidth}}{
        \cellcolor[rgb]{0.9,0.9,0.9}{
            \makecell[{{p{\linewidth}}}]{
                \texttt{\tiny{[GM$\langle$P2]}}
                \texttt{ВІДГАДКА: міст, бар, лікарня} \\
            }
        }
    }
     \\ \\

    \theutterance \stepcounter{utterance}  
    & & & \multicolumn{2}{p{0.3\linewidth}}{
        \cellcolor[rgb]{0.9,0.9,0.9}{
            \makecell[{{p{\linewidth}}}]{
                \texttt{\tiny{[GM$|$GM]}}
                ['міст', 'бар', 'лікарня']
            }
        }
    }
    & & \\ \\

    \theutterance \stepcounter{utterance}  
    & & & \multicolumn{2}{p{0.3\linewidth}}{
        \cellcolor[rgb]{0.9,0.9,0.9}{
            \makecell[{{p{\linewidth}}}]{
                \texttt{\tiny{[GM$|$GM]}}
                {'word': 'міст', 'assignment': 'team'}
            }
        }
    }
    & & \\ \\

    \theutterance \stepcounter{utterance}  
    & & & \multicolumn{2}{p{0.3\linewidth}}{
        \cellcolor[rgb]{0.9,0.9,0.9}{
            \makecell[{{p{\linewidth}}}]{
                \texttt{\tiny{[GM$|$GM]}}
                {'word': 'міст', 'assignment': 'team'}
            }
        }
    }
    & & \\ \\

    \theutterance \stepcounter{utterance}  
    & & & \multicolumn{2}{p{0.3\linewidth}}{
        \cellcolor[rgb]{0.9,0.9,0.9}{
            \makecell[{{p{\linewidth}}}]{
                \texttt{\tiny{[GM$|$GM]}}
                {'word': 'бар', 'assignment': 'innocent'}
            }
        }
    }
    & & \\ \\

    \theutterance \stepcounter{utterance}  
    & & & \multicolumn{2}{p{0.3\linewidth}}{
        \cellcolor[rgb]{0.9,0.9,0.9}{
            \makecell[{{p{\linewidth}}}]{
                \texttt{\tiny{[GM$|$GM]}}
                \texttt{бар} \\
            }
        }
    }
    & & \\ \\

    \theutterance \stepcounter{utterance}  
    & & & \multicolumn{2}{p{0.3\linewidth}}{
        \cellcolor[rgb]{0.9,0.9,0.9}{
            \makecell[{{p{\linewidth}}}]{
                \texttt{\tiny{[GM$|$GM]}}
                {'word': 'лимон', 'assignment': 'opponent'}
            }
        }
    }
    & & \\ \\

    \theutterance \stepcounter{utterance}  
    & & & \multicolumn{2}{p{0.3\linewidth}}{
        \cellcolor[rgb]{0.9,0.9,0.9}{
            \makecell[{{p{\linewidth}}}]{
                \texttt{\tiny{[GM$|$GM]}}
                {'hidden': {'team': ['революція', 'дошка', 'курча', 'відьма', 'сила', 'клен', 'контракт', 'ліжко'], 'innocent': ['лід', 'лепрекон', 'екран', 'лікарня', 'москва', 'вентилятор', 'машина'], 'opponent': ['злодій', 'залізо', 'прибулець', 'концерт', 'басейн', 'африка', 'стопа'], 'assassin': []}, 'revealed': {'team': {'team': ['міст'], 'innocent': ['бар'], 'opponent': [], 'assassin': []}, 'opponent': {'team': [], 'innocent': [], 'opponent': ['лимон'], 'assassin': []}}}
            }
        }
    }
    & & \\ \\

    \theutterance \stepcounter{utterance}  
    & & \multicolumn{4}{p{0.6\linewidth}}{
        \cellcolor[rgb]{0.9,0.9,0.9}{
            \makecell[{{p{\linewidth}}}]{
                \texttt{\tiny{[P1$\langle$GM]}}
                \texttt{Слово міст було відгадано правильно. Слово бар було відгадано, але це нейтральне слово. Хід вашого товариша по команді закінчився.} \\
\texttt{Тепер надайте ще одну підказку, пов'язану з деякими з ваших командних слів, що залишилися, та список пов'язаних цільових слів. Пам'ятайте: почніть підказку з 'ПІДКАЗКА: ', перейдіть на новий рядок і почніть ваш список цільових слів, розділених комами, з 'ЦІЛІ: '. Зверніть увагу: деякі слова було вилучено зі списків порівняно з попередніми запитами.} \\
\\ 
\texttt{Ваші командні слова, що залишилися: революція, дошка, курча, відьма, сила, клен, контракт, ліжко.} \\
\texttt{Слова, що залишилися для вашого суперника: злодій, лимон, залізо, прибулець, концерт, басейн, африка, стопа.} \\
\texttt{Слова{-}відволікачі, що залишилися: лід, лепрекон, екран, лікарня, москва, вентилятор, машина.} \\
\texttt{Слова{-}вбивці, що залишилися: .} \\
            }
        }
    }
    & \\ \\

    \theutterance \stepcounter{utterance}  
    & \multicolumn{4}{p{0.6\linewidth}}{
        \cellcolor[rgb]{0.9,0.9,0.9}{
            \makecell[{{p{\linewidth}}}]{
                \texttt{\tiny{[P1$\rangle$GM]}}
                \texttt{ПІДКАЗКА: дерево} \\
\texttt{ЦІЛІ: клен} \\
            }
        }
    }
    & & \\ \\

    \theutterance \stepcounter{utterance}  
    & & & \multicolumn{2}{p{0.3\linewidth}}{
        \cellcolor[rgb]{0.9,0.9,0.9}{
            \makecell[{{p{\linewidth}}}]{
                \texttt{\tiny{[GM$|$GM]}}
                \texttt{дерево} \\
            }
        }
    }
    & & \\ \\

    \theutterance \stepcounter{utterance}  
    & & & \multicolumn{2}{p{0.3\linewidth}}{
        \cellcolor[rgb]{0.9,0.9,0.9}{
            \makecell[{{p{\linewidth}}}]{
                \texttt{\tiny{[GM$|$GM]}}
                ['клен']
            }
        }
    }
    & & \\ \\

    \theutterance \stepcounter{utterance}  
    & & & \multicolumn{2}{p{0.3\linewidth}}{
        \cellcolor[rgb]{0.9,0.9,0.9}{
            \makecell[{{p{\linewidth}}}]{
                \texttt{\tiny{[GM$|$GM]}}
                {'word': 'клен', 'assignment': 'team'}
            }
        }
    }
    & & \\ \\

    \theutterance \stepcounter{utterance}  
    & & \multicolumn{4}{p{0.6\linewidth}}{
        \cellcolor[rgb]{0.9,0.9,0.9}{
            \makecell[{{p{\linewidth}}}]{
                \texttt{\tiny{[GM$\rangle$P2]}}
                \texttt{Слово міст було відгадано правильно. Слово бар було відгадано, але це нейтральне слово. Ваш хід закінчився.} \\
\texttt{Тепер надайте ще один список слів, розділених комами, щонайменше з 1 і щонайбільше з 1 слів із наведеного нижче списку слів, які найкраще пов'язані або найбільш тісно асоціюються зі словом 'дерево'. Пам'ятайте, що потрібно почати вашу відповідь з ВІДГАДКА: . Зауваження: деякі слова було вилучено зі списку порівняно з попередніми запитами.} \\
\\ 
\texttt{ліжко, сила, відьма, клен, стопа, машина, лепрекон, басейн, революція, лікарня, залізо, концерт, курча, лід, злодій, вентилятор, контракт, африка, москва, екран, дошка, прибулець} \\
            }
        }
    }
    & \\ \\

    \theutterance \stepcounter{utterance}  
    & & & \multicolumn{4}{p{0.6\linewidth}}{
        \cellcolor[rgb]{0.9,0.9,0.9}{
            \makecell[{{p{\linewidth}}}]{
                \texttt{\tiny{[GM$\langle$P2]}}
                \texttt{ВІДГАДКА: клен} \\
            }
        }
    }
     \\ \\

    \theutterance \stepcounter{utterance}  
    & & & \multicolumn{2}{p{0.3\linewidth}}{
        \cellcolor[rgb]{0.9,0.9,0.9}{
            \makecell[{{p{\linewidth}}}]{
                \texttt{\tiny{[GM$|$GM]}}
                ['клен']
            }
        }
    }
    & & \\ \\

    \theutterance \stepcounter{utterance}  
    & & & \multicolumn{2}{p{0.3\linewidth}}{
        \cellcolor[rgb]{0.9,0.9,0.9}{
            \makecell[{{p{\linewidth}}}]{
                \texttt{\tiny{[GM$|$GM]}}
                {'word': 'клен', 'assignment': 'team'}
            }
        }
    }
    & & \\ \\

    \theutterance \stepcounter{utterance}  
    & & & \multicolumn{2}{p{0.3\linewidth}}{
        \cellcolor[rgb]{0.9,0.9,0.9}{
            \makecell[{{p{\linewidth}}}]{
                \texttt{\tiny{[GM$|$GM]}}
                {'word': 'клен', 'assignment': 'team'}
            }
        }
    }
    & & \\ \\

    \theutterance \stepcounter{utterance}  
    & & & \multicolumn{2}{p{0.3\linewidth}}{
        \cellcolor[rgb]{0.9,0.9,0.9}{
            \makecell[{{p{\linewidth}}}]{
                \texttt{\tiny{[GM$|$GM]}}
                {'word': 'злодій', 'assignment': 'opponent'}
            }
        }
    }
    & & \\ \\

    \theutterance \stepcounter{utterance}  
    & & & \multicolumn{2}{p{0.3\linewidth}}{
        \cellcolor[rgb]{0.9,0.9,0.9}{
            \makecell[{{p{\linewidth}}}]{
                \texttt{\tiny{[GM$|$GM]}}
                {'hidden': {'team': ['революція', 'дошка', 'курча', 'відьма', 'сила', 'контракт', 'ліжко'], 'innocent': ['лід', 'лепрекон', 'екран', 'лікарня', 'москва', 'вентилятор', 'машина'], 'opponent': ['залізо', 'прибулець', 'концерт', 'басейн', 'африка', 'стопа'], 'assassin': []}, 'revealed': {'team': {'team': ['міст', 'клен'], 'innocent': ['бар'], 'opponent': [], 'assassin': []}, 'opponent': {'team': [], 'innocent': [], 'opponent': ['лимон', 'злодій'], 'assassin': []}}}
            }
        }
    }
    & & \\ \\

    \theutterance \stepcounter{utterance}  
    & & \multicolumn{4}{p{0.6\linewidth}}{
        \cellcolor[rgb]{0.9,0.9,0.9}{
            \makecell[{{p{\linewidth}}}]{
                \texttt{\tiny{[P1$\langle$GM]}}
                \texttt{Слово клен було відгадано правильно. Хід вашого товариша по команді закінчився.} \\
\texttt{Тепер надайте ще одну підказку, пов'язану з деякими з ваших командних слів, що залишилися, та список пов'язаних цільових слів. Пам'ятайте: почніть підказку з 'ПІДКАЗКА: ', перейдіть на новий рядок і почніть ваш список цільових слів, розділених комами, з 'ЦІЛІ: '. Зверніть увагу: деякі слова було вилучено зі списків порівняно з попередніми запитами.} \\
\\ 
\texttt{Ваші командні слова, що залишилися: революція, дошка, курча, відьма, сила, контракт, ліжко.} \\
\texttt{Слова, що залишилися для вашого суперника: злодій, залізо, прибулець, концерт, басейн, африка, стопа.} \\
\texttt{Слова{-}відволікачі, що залишилися: лід, лепрекон, екран, лікарня, москва, вентилятор, машина.} \\
\texttt{Слова{-}вбивці, що залишилися: .} \\
            }
        }
    }
    & \\ \\

    \theutterance \stepcounter{utterance}  
    & \multicolumn{4}{p{0.6\linewidth}}{
        \cellcolor[rgb]{0.9,0.9,0.9}{
            \makecell[{{p{\linewidth}}}]{
                \texttt{\tiny{[P1$\rangle$GM]}}
                \texttt{ПІДКАЗКА: договір} \\
\texttt{ЦІЛІ: контракт} \\
            }
        }
    }
    & & \\ \\

    \theutterance \stepcounter{utterance}  
    & & & \multicolumn{2}{p{0.3\linewidth}}{
        \cellcolor[rgb]{0.9,0.9,0.9}{
            \makecell[{{p{\linewidth}}}]{
                \texttt{\tiny{[GM$|$GM]}}
                \texttt{договір} \\
            }
        }
    }
    & & \\ \\

    \theutterance \stepcounter{utterance}  
    & & & \multicolumn{2}{p{0.3\linewidth}}{
        \cellcolor[rgb]{0.9,0.9,0.9}{
            \makecell[{{p{\linewidth}}}]{
                \texttt{\tiny{[GM$|$GM]}}
                ['контракт']
            }
        }
    }
    & & \\ \\

    \theutterance \stepcounter{utterance}  
    & & & \multicolumn{2}{p{0.3\linewidth}}{
        \cellcolor[rgb]{0.9,0.9,0.9}{
            \makecell[{{p{\linewidth}}}]{
                \texttt{\tiny{[GM$|$GM]}}
                {'word': 'контракт', 'assignment': 'team'}
            }
        }
    }
    & & \\ \\

    \theutterance \stepcounter{utterance}  
    & & \multicolumn{4}{p{0.6\linewidth}}{
        \cellcolor[rgb]{0.9,0.9,0.9}{
            \makecell[{{p{\linewidth}}}]{
                \texttt{\tiny{[GM$\rangle$P2]}}
                \texttt{Слово клен було відгадано правильно. Ваш хід закінчився.} \\
\texttt{Тепер надайте ще один список слів, розділених комами, щонайменше з 1 і щонайбільше з 1 слів із наведеного нижче списку слів, які найкраще пов'язані або найбільш тісно асоціюються зі словом 'договір'. Пам'ятайте, що потрібно почати вашу відповідь з ВІДГАДКА: . Зауваження: деякі слова було вилучено зі списку порівняно з попередніми запитами.} \\
\\ 
\texttt{ліжко, сила, відьма, стопа, машина, лепрекон, басейн, революція, лікарня, залізо, концерт, курча, лід, вентилятор, контракт, африка, москва, екран, дошка, прибулець} \\
            }
        }
    }
    & \\ \\

    \theutterance \stepcounter{utterance}  
    & & & \multicolumn{4}{p{0.6\linewidth}}{
        \cellcolor[rgb]{0.9,0.9,0.9}{
            \makecell[{{p{\linewidth}}}]{
                \texttt{\tiny{[GM$\langle$P2]}}
                \texttt{ВІДГАДКА: контракт} \\
            }
        }
    }
     \\ \\

    \theutterance \stepcounter{utterance}  
    & & & \multicolumn{2}{p{0.3\linewidth}}{
        \cellcolor[rgb]{0.9,0.9,0.9}{
            \makecell[{{p{\linewidth}}}]{
                \texttt{\tiny{[GM$|$GM]}}
                ['контракт']
            }
        }
    }
    & & \\ \\

    \theutterance \stepcounter{utterance}  
    & & & \multicolumn{2}{p{0.3\linewidth}}{
        \cellcolor[rgb]{0.9,0.9,0.9}{
            \makecell[{{p{\linewidth}}}]{
                \texttt{\tiny{[GM$|$GM]}}
                {'word': 'контракт', 'assignment': 'team'}
            }
        }
    }
    & & \\ \\

    \theutterance \stepcounter{utterance}  
    & & & \multicolumn{2}{p{0.3\linewidth}}{
        \cellcolor[rgb]{0.9,0.9,0.9}{
            \makecell[{{p{\linewidth}}}]{
                \texttt{\tiny{[GM$|$GM]}}
                {'word': 'контракт', 'assignment': 'team'}
            }
        }
    }
    & & \\ \\

    \theutterance \stepcounter{utterance}  
    & & & \multicolumn{2}{p{0.3\linewidth}}{
        \cellcolor[rgb]{0.9,0.9,0.9}{
            \makecell[{{p{\linewidth}}}]{
                \texttt{\tiny{[GM$|$GM]}}
                {'word': 'стопа', 'assignment': 'opponent'}
            }
        }
    }
    & & \\ \\

    \theutterance \stepcounter{utterance}  
    & & & \multicolumn{2}{p{0.3\linewidth}}{
        \cellcolor[rgb]{0.9,0.9,0.9}{
            \makecell[{{p{\linewidth}}}]{
                \texttt{\tiny{[GM$|$GM]}}
                {'hidden': {'team': ['революція', 'дошка', 'курча', 'відьма', 'сила', 'ліжко'], 'innocent': ['лід', 'лепрекон', 'екран', 'лікарня', 'москва', 'вентилятор', 'машина'], 'opponent': ['залізо', 'прибулець', 'концерт', 'басейн', 'африка'], 'assassin': []}, 'revealed': {'team': {'team': ['міст', 'клен', 'контракт'], 'innocent': ['бар'], 'opponent': [], 'assassin': []}, 'opponent': {'team': [], 'innocent': [], 'opponent': ['лимон', 'злодій', 'стопа'], 'assassin': []}}}
            }
        }
    }
    & & \\ \\

    \theutterance \stepcounter{utterance}  
    & & \multicolumn{4}{p{0.6\linewidth}}{
        \cellcolor[rgb]{0.9,0.9,0.9}{
            \makecell[{{p{\linewidth}}}]{
                \texttt{\tiny{[P1$\langle$GM]}}
                \texttt{Слово контракт було відгадано правильно. Хід вашого товариша по команді закінчився.} \\
\texttt{Тепер надайте ще одну підказку, пов'язану з деякими з ваших командних слів, що залишилися, та список пов'язаних цільових слів. Пам'ятайте: почніть підказку з 'ПІДКАЗКА: ', перейдіть на новий рядок і почніть ваш список цільових слів, розділених комами, з 'ЦІЛІ: '. Зверніть увагу: деякі слова було вилучено зі списків порівняно з попередніми запитами.} \\
\\ 
\texttt{Ваші командні слова, що залишилися: революція, дошка, курча, відьма, сила, ліжко.} \\
\texttt{Слова, що залишилися для вашого суперника: залізо, прибулець, концерт, басейн, африка, стопа.} \\
\texttt{Слова{-}відволікачі, що залишилися: лід, лепрекон, екран, лікарня, москва, вентилятор, машина.} \\
\texttt{Слова{-}вбивці, що залишилися: .} \\
            }
        }
    }
    & \\ \\

    \theutterance \stepcounter{utterance}  
    & \multicolumn{4}{p{0.6\linewidth}}{
        \cellcolor[rgb]{0.9,0.9,0.9}{
            \makecell[{{p{\linewidth}}}]{
                \texttt{\tiny{[P1$\rangle$GM]}}
                \texttt{ПІДКАЗКА: спальня} \\
\texttt{ЦІЛІ: ліжко} \\
            }
        }
    }
    & & \\ \\

    \theutterance \stepcounter{utterance}  
    & & & \multicolumn{2}{p{0.3\linewidth}}{
        \cellcolor[rgb]{0.9,0.9,0.9}{
            \makecell[{{p{\linewidth}}}]{
                \texttt{\tiny{[GM$|$GM]}}
                \texttt{спальня} \\
            }
        }
    }
    & & \\ \\

    \theutterance \stepcounter{utterance}  
    & & & \multicolumn{2}{p{0.3\linewidth}}{
        \cellcolor[rgb]{0.9,0.9,0.9}{
            \makecell[{{p{\linewidth}}}]{
                \texttt{\tiny{[GM$|$GM]}}
                ['ліжко']
            }
        }
    }
    & & \\ \\

    \theutterance \stepcounter{utterance}  
    & & & \multicolumn{2}{p{0.3\linewidth}}{
        \cellcolor[rgb]{0.9,0.9,0.9}{
            \makecell[{{p{\linewidth}}}]{
                \texttt{\tiny{[GM$|$GM]}}
                {'word': 'ліжко', 'assignment': 'team'}
            }
        }
    }
    & & \\ \\

    \theutterance \stepcounter{utterance}  
    & & \multicolumn{4}{p{0.6\linewidth}}{
        \cellcolor[rgb]{0.9,0.9,0.9}{
            \makecell[{{p{\linewidth}}}]{
                \texttt{\tiny{[GM$\rangle$P2]}}
                \texttt{Слово контракт було відгадано правильно. Ваш хід закінчився.} \\
\texttt{Тепер надайте ще один список слів, розділених комами, щонайменше з 1 і щонайбільше з 1 слів із наведеного нижче списку слів, які найкраще пов'язані або найбільш тісно асоціюються зі словом 'спальня'. Пам'ятайте, що потрібно почати вашу відповідь з ВІДГАДКА: . Зауваження: деякі слова було вилучено зі списку порівняно з попередніми запитами.} \\
\\ 
\texttt{ліжко, сила, відьма, машина, лепрекон, басейн, революція, лікарня, залізо, концерт, курча, лід, вентилятор, африка, москва, екран, дошка, прибулець} \\
            }
        }
    }
    & \\ \\

    \theutterance \stepcounter{utterance}  
    & & & \multicolumn{4}{p{0.6\linewidth}}{
        \cellcolor[rgb]{0.9,0.9,0.9}{
            \makecell[{{p{\linewidth}}}]{
                \texttt{\tiny{[GM$\langle$P2]}}
                \texttt{ВІДГАДКА: ліжко} \\
            }
        }
    }
     \\ \\

    \theutterance \stepcounter{utterance}  
    & & & \multicolumn{2}{p{0.3\linewidth}}{
        \cellcolor[rgb]{0.9,0.9,0.9}{
            \makecell[{{p{\linewidth}}}]{
                \texttt{\tiny{[GM$|$GM]}}
                ['ліжко']
            }
        }
    }
    & & \\ \\

    \theutterance \stepcounter{utterance}  
    & & & \multicolumn{2}{p{0.3\linewidth}}{
        \cellcolor[rgb]{0.9,0.9,0.9}{
            \makecell[{{p{\linewidth}}}]{
                \texttt{\tiny{[GM$|$GM]}}
                {'word': 'ліжко', 'assignment': 'team'}
            }
        }
    }
    & & \\ \\

    \theutterance \stepcounter{utterance}  
    & & & \multicolumn{2}{p{0.3\linewidth}}{
        \cellcolor[rgb]{0.9,0.9,0.9}{
            \makecell[{{p{\linewidth}}}]{
                \texttt{\tiny{[GM$|$GM]}}
                {'word': 'ліжко', 'assignment': 'team'}
            }
        }
    }
    & & \\ \\

    \theutterance \stepcounter{utterance}  
    & & & \multicolumn{2}{p{0.3\linewidth}}{
        \cellcolor[rgb]{0.9,0.9,0.9}{
            \makecell[{{p{\linewidth}}}]{
                \texttt{\tiny{[GM$|$GM]}}
                {'word': 'концерт', 'assignment': 'opponent'}
            }
        }
    }
    & & \\ \\

    \theutterance \stepcounter{utterance}  
    & & & \multicolumn{2}{p{0.3\linewidth}}{
        \cellcolor[rgb]{0.9,0.9,0.9}{
            \makecell[{{p{\linewidth}}}]{
                \texttt{\tiny{[GM$|$GM]}}
                {'hidden': {'team': ['революція', 'дошка', 'курча', 'відьма', 'сила'], 'innocent': ['лід', 'лепрекон', 'екран', 'лікарня', 'москва', 'вентилятор', 'машина'], 'opponent': ['залізо', 'прибулець', 'басейн', 'африка'], 'assassin': []}, 'revealed': {'team': {'team': ['міст', 'клен', 'контракт', 'ліжко'], 'innocent': ['бар'], 'opponent': [], 'assassin': []}, 'opponent': {'team': [], 'innocent': [], 'opponent': ['лимон', 'злодій', 'стопа', 'концерт'], 'assassin': []}}}
            }
        }
    }
    & & \\ \\

    \theutterance \stepcounter{utterance}  
    & & \multicolumn{4}{p{0.6\linewidth}}{
        \cellcolor[rgb]{0.9,0.9,0.9}{
            \makecell[{{p{\linewidth}}}]{
                \texttt{\tiny{[P1$\langle$GM]}}
                \texttt{Слово ліжко було відгадано правильно. Хід вашого товариша по команді закінчився.} \\
\texttt{Тепер надайте ще одну підказку, пов'язану з деякими з ваших командних слів, що залишилися, та список пов'язаних цільових слів. Пам'ятайте: почніть підказку з 'ПІДКАЗКА: ', перейдіть на новий рядок і почніть ваш список цільових слів, розділених комами, з 'ЦІЛІ: '. Зверніть увагу: деякі слова було вилучено зі списків порівняно з попередніми запитами.} \\
\\ 
\texttt{Ваші командні слова, що залишилися: революція, дошка, курча, відьма, сила.} \\
\texttt{Слова, що залишилися для вашого суперника: залізо, прибулець, концерт, басейн, африка.} \\
\texttt{Слова{-}відволікачі, що залишилися: лід, лепрекон, екран, лікарня, москва, вентилятор, машина.} \\
\texttt{Слова{-}вбивці, що залишилися: .} \\
            }
        }
    }
    & \\ \\

    \theutterance \stepcounter{utterance}  
    & \multicolumn{4}{p{0.6\linewidth}}{
        \cellcolor[rgb]{0.9,0.9,0.9}{
            \makecell[{{p{\linewidth}}}]{
                \texttt{\tiny{[P1$\rangle$GM]}}
                \texttt{ПІДКАЗКА: магія} \\
\texttt{ЦІЛІ: відьма} \\
            }
        }
    }
    & & \\ \\

    \theutterance \stepcounter{utterance}  
    & & & \multicolumn{2}{p{0.3\linewidth}}{
        \cellcolor[rgb]{0.9,0.9,0.9}{
            \makecell[{{p{\linewidth}}}]{
                \texttt{\tiny{[GM$|$GM]}}
                \texttt{магія} \\
            }
        }
    }
    & & \\ \\

    \theutterance \stepcounter{utterance}  
    & & & \multicolumn{2}{p{0.3\linewidth}}{
        \cellcolor[rgb]{0.9,0.9,0.9}{
            \makecell[{{p{\linewidth}}}]{
                \texttt{\tiny{[GM$|$GM]}}
                ['відьма']
            }
        }
    }
    & & \\ \\

    \theutterance \stepcounter{utterance}  
    & & & \multicolumn{2}{p{0.3\linewidth}}{
        \cellcolor[rgb]{0.9,0.9,0.9}{
            \makecell[{{p{\linewidth}}}]{
                \texttt{\tiny{[GM$|$GM]}}
                {'word': 'відьма', 'assignment': 'team'}
            }
        }
    }
    & & \\ \\

    \theutterance \stepcounter{utterance}  
    & & \multicolumn{4}{p{0.6\linewidth}}{
        \cellcolor[rgb]{0.9,0.9,0.9}{
            \makecell[{{p{\linewidth}}}]{
                \texttt{\tiny{[GM$\rangle$P2]}}
                \texttt{Слово ліжко було відгадано правильно. Ваш хід закінчився.} \\
\texttt{Тепер надайте ще один список слів, розділених комами, щонайменше з 1 і щонайбільше з 1 слів із наведеного нижче списку слів, які найкраще пов'язані або найбільш тісно асоціюються зі словом 'магія'. Пам'ятайте, що потрібно почати вашу відповідь з ВІДГАДКА: . Зауваження: деякі слова було вилучено зі списку порівняно з попередніми запитами.} \\
\\ 
\texttt{сила, відьма, машина, лепрекон, басейн, революція, лікарня, залізо, курча, лід, вентилятор, африка, москва, екран, дошка, прибулець} \\
            }
        }
    }
    & \\ \\

    \theutterance \stepcounter{utterance}  
    & & & \multicolumn{4}{p{0.6\linewidth}}{
        \cellcolor[rgb]{0.9,0.9,0.9}{
            \makecell[{{p{\linewidth}}}]{
                \texttt{\tiny{[GM$\langle$P2]}}
                \texttt{ВІДГАДКА: відьма} \\
            }
        }
    }
     \\ \\

    \theutterance \stepcounter{utterance}  
    & & & \multicolumn{2}{p{0.3\linewidth}}{
        \cellcolor[rgb]{0.9,0.9,0.9}{
            \makecell[{{p{\linewidth}}}]{
                \texttt{\tiny{[GM$|$GM]}}
                ['відьма']
            }
        }
    }
    & & \\ \\

    \theutterance \stepcounter{utterance}  
    & & & \multicolumn{2}{p{0.3\linewidth}}{
        \cellcolor[rgb]{0.9,0.9,0.9}{
            \makecell[{{p{\linewidth}}}]{
                \texttt{\tiny{[GM$|$GM]}}
                {'word': 'відьма', 'assignment': 'team'}
            }
        }
    }
    & & \\ \\

    \theutterance \stepcounter{utterance}  
    & & & \multicolumn{2}{p{0.3\linewidth}}{
        \cellcolor[rgb]{0.9,0.9,0.9}{
            \makecell[{{p{\linewidth}}}]{
                \texttt{\tiny{[GM$|$GM]}}
                {'word': 'відьма', 'assignment': 'team'}
            }
        }
    }
    & & \\ \\

    \theutterance \stepcounter{utterance}  
    & & & \multicolumn{2}{p{0.3\linewidth}}{
        \cellcolor[rgb]{0.9,0.9,0.9}{
            \makecell[{{p{\linewidth}}}]{
                \texttt{\tiny{[GM$|$GM]}}
                {'word': 'прибулець', 'assignment': 'opponent'}
            }
        }
    }
    & & \\ \\

    \theutterance \stepcounter{utterance}  
    & & & \multicolumn{2}{p{0.3\linewidth}}{
        \cellcolor[rgb]{0.9,0.9,0.9}{
            \makecell[{{p{\linewidth}}}]{
                \texttt{\tiny{[GM$|$GM]}}
                {'hidden': {'team': ['революція', 'дошка', 'курча', 'сила'], 'innocent': ['лід', 'лепрекон', 'екран', 'лікарня', 'москва', 'вентилятор', 'машина'], 'opponent': ['залізо', 'басейн', 'африка'], 'assassin': []}, 'revealed': {'team': {'team': ['міст', 'клен', 'контракт', 'ліжко', 'відьма'], 'innocent': ['бар'], 'opponent': [], 'assassin': []}, 'opponent': {'team': [], 'innocent': [], 'opponent': ['лимон', 'злодій', 'стопа', 'концерт', 'прибулець'], 'assassin': []}}}
            }
        }
    }
    & & \\ \\

    \theutterance \stepcounter{utterance}  
    & & \multicolumn{4}{p{0.6\linewidth}}{
        \cellcolor[rgb]{0.9,0.9,0.9}{
            \makecell[{{p{\linewidth}}}]{
                \texttt{\tiny{[P1$\langle$GM]}}
                \texttt{Слово відьма було відгадано правильно. Хід вашого товариша по команді закінчився.} \\
\texttt{Тепер надайте ще одну підказку, пов'язану з деякими з ваших командних слів, що залишилися, та список пов'язаних цільових слів. Пам'ятайте: почніть підказку з 'ПІДКАЗКА: ', перейдіть на новий рядок і почніть ваш список цільових слів, розділених комами, з 'ЦІЛІ: '. Зверніть увагу: деякі слова було вилучено зі списків порівняно з попередніми запитами.} \\
\\ 
\texttt{Ваші командні слова, що залишилися: революція, дошка, курча, сила.} \\
\texttt{Слова, що залишилися для вашого суперника: залізо, прибулець, басейн, африка.} \\
\texttt{Слова{-}відволікачі, що залишилися: лід, лепрекон, екран, лікарня, москва, вентилятор, машина.} \\
\texttt{Слова{-}вбивці, що залишилися: .} \\
            }
        }
    }
    & \\ \\

    \theutterance \stepcounter{utterance}  
    & \multicolumn{4}{p{0.6\linewidth}}{
        \cellcolor[rgb]{0.9,0.9,0.9}{
            \makecell[{{p{\linewidth}}}]{
                \texttt{\tiny{[P1$\rangle$GM]}}
                \texttt{ПІДКАЗКА: їжа} \\
\texttt{ЦІЛІ: курча} \\
            }
        }
    }
    & & \\ \\

    \theutterance \stepcounter{utterance}  
    & & & \multicolumn{2}{p{0.3\linewidth}}{
        \cellcolor[rgb]{0.9,0.9,0.9}{
            \makecell[{{p{\linewidth}}}]{
                \texttt{\tiny{[GM$|$GM]}}
                \texttt{їжа} \\
            }
        }
    }
    & & \\ \\

    \theutterance \stepcounter{utterance}  
    & & & \multicolumn{2}{p{0.3\linewidth}}{
        \cellcolor[rgb]{0.9,0.9,0.9}{
            \makecell[{{p{\linewidth}}}]{
                \texttt{\tiny{[GM$|$GM]}}
                ['курча']
            }
        }
    }
    & & \\ \\

    \theutterance \stepcounter{utterance}  
    & & & \multicolumn{2}{p{0.3\linewidth}}{
        \cellcolor[rgb]{0.9,0.9,0.9}{
            \makecell[{{p{\linewidth}}}]{
                \texttt{\tiny{[GM$|$GM]}}
                {'word': 'курча', 'assignment': 'team'}
            }
        }
    }
    & & \\ \\

    \theutterance \stepcounter{utterance}  
    & & \multicolumn{4}{p{0.6\linewidth}}{
        \cellcolor[rgb]{0.9,0.9,0.9}{
            \makecell[{{p{\linewidth}}}]{
                \texttt{\tiny{[GM$\rangle$P2]}}
                \texttt{Слово відьма було відгадано правильно. Ваш хід закінчився.} \\
\texttt{Тепер надайте ще один список слів, розділених комами, щонайменше з 1 і щонайбільше з 1 слів із наведеного нижче списку слів, які найкраще пов'язані або найбільш тісно асоціюються зі словом 'їжа'. Пам'ятайте, що потрібно почати вашу відповідь з ВІДГАДКА: . Зауваження: деякі слова було вилучено зі списку порівняно з попередніми запитами.} \\
\\ 
\texttt{сила, машина, лепрекон, басейн, революція, лікарня, залізо, курча, лід, вентилятор, африка, москва, екран, дошка} \\
            }
        }
    }
    & \\ \\

    \theutterance \stepcounter{utterance}  
    & & & \multicolumn{4}{p{0.6\linewidth}}{
        \cellcolor[rgb]{0.9,0.9,0.9}{
            \makecell[{{p{\linewidth}}}]{
                \texttt{\tiny{[GM$\langle$P2]}}
                \texttt{ВІДГАДКА: курча} \\
            }
        }
    }
     \\ \\

    \theutterance \stepcounter{utterance}  
    & & & \multicolumn{2}{p{0.3\linewidth}}{
        \cellcolor[rgb]{0.9,0.9,0.9}{
            \makecell[{{p{\linewidth}}}]{
                \texttt{\tiny{[GM$|$GM]}}
                ['курча']
            }
        }
    }
    & & \\ \\

    \theutterance \stepcounter{utterance}  
    & & & \multicolumn{2}{p{0.3\linewidth}}{
        \cellcolor[rgb]{0.9,0.9,0.9}{
            \makecell[{{p{\linewidth}}}]{
                \texttt{\tiny{[GM$|$GM]}}
                {'word': 'курча', 'assignment': 'team'}
            }
        }
    }
    & & \\ \\

    \theutterance \stepcounter{utterance}  
    & & & \multicolumn{2}{p{0.3\linewidth}}{
        \cellcolor[rgb]{0.9,0.9,0.9}{
            \makecell[{{p{\linewidth}}}]{
                \texttt{\tiny{[GM$|$GM]}}
                {'word': 'курча', 'assignment': 'team'}
            }
        }
    }
    & & \\ \\

    \theutterance \stepcounter{utterance}  
    & & & \multicolumn{2}{p{0.3\linewidth}}{
        \cellcolor[rgb]{0.9,0.9,0.9}{
            \makecell[{{p{\linewidth}}}]{
                \texttt{\tiny{[GM$|$GM]}}
                {'word': 'залізо', 'assignment': 'opponent'}
            }
        }
    }
    & & \\ \\

    \theutterance \stepcounter{utterance}  
    & & & \multicolumn{2}{p{0.3\linewidth}}{
        \cellcolor[rgb]{0.9,0.9,0.9}{
            \makecell[{{p{\linewidth}}}]{
                \texttt{\tiny{[GM$|$GM]}}
                {'hidden': {'team': ['революція', 'дошка', 'сила'], 'innocent': ['лід', 'лепрекон', 'екран', 'лікарня', 'москва', 'вентилятор', 'машина'], 'opponent': ['басейн', 'африка'], 'assassin': []}, 'revealed': {'team': {'team': ['міст', 'клен', 'контракт', 'ліжко', 'відьма', 'курча'], 'innocent': ['бар'], 'opponent': [], 'assassin': []}, 'opponent': {'team': [], 'innocent': [], 'opponent': ['лимон', 'злодій', 'стопа', 'концерт', 'прибулець', 'залізо'], 'assassin': []}}}
            }
        }
    }
    & & \\ \\

    \theutterance \stepcounter{utterance}  
    & & \multicolumn{4}{p{0.6\linewidth}}{
        \cellcolor[rgb]{0.9,0.9,0.9}{
            \makecell[{{p{\linewidth}}}]{
                \texttt{\tiny{[P1$\langle$GM]}}
                \texttt{Слово курча було відгадано правильно. Хід вашого товариша по команді закінчився.} \\
\texttt{Тепер надайте ще одну підказку, пов'язану з деякими з ваших командних слів, що залишилися, та список пов'язаних цільових слів. Пам'ятайте: почніть підказку з 'ПІДКАЗКА: ', перейдіть на новий рядок і почніть ваш список цільових слів, розділених комами, з 'ЦІЛІ: '. Зверніть увагу: деякі слова було вилучено зі списків порівняно з попередніми запитами.} \\
\\ 
\texttt{Ваші командні слова, що залишилися: революція, дошка, сила.} \\
\texttt{Слова, що залишилися для вашого суперника: залізо, басейн, африка.} \\
\texttt{Слова{-}відволікачі, що залишилися: лід, лепрекон, екран, лікарня, москва, вентилятор, машина.} \\
\texttt{Слова{-}вбивці, що залишилися: .} \\
            }
        }
    }
    & \\ \\

    \theutterance \stepcounter{utterance}  
    & \multicolumn{4}{p{0.6\linewidth}}{
        \cellcolor[rgb]{0.9,0.9,0.9}{
            \makecell[{{p{\linewidth}}}]{
                \texttt{\tiny{[P1$\rangle$GM]}}
                \texttt{ПІДКАЗКА: спорт} \\
\texttt{ЦІЛІ: сила} \\
            }
        }
    }
    & & \\ \\

    \theutterance \stepcounter{utterance}  
    & & & \multicolumn{2}{p{0.3\linewidth}}{
        \cellcolor[rgb]{0.9,0.9,0.9}{
            \makecell[{{p{\linewidth}}}]{
                \texttt{\tiny{[GM$|$GM]}}
                \texttt{спорт} \\
            }
        }
    }
    & & \\ \\

    \theutterance \stepcounter{utterance}  
    & & & \multicolumn{2}{p{0.3\linewidth}}{
        \cellcolor[rgb]{0.9,0.9,0.9}{
            \makecell[{{p{\linewidth}}}]{
                \texttt{\tiny{[GM$|$GM]}}
                ['сила']
            }
        }
    }
    & & \\ \\

    \theutterance \stepcounter{utterance}  
    & & & \multicolumn{2}{p{0.3\linewidth}}{
        \cellcolor[rgb]{0.9,0.9,0.9}{
            \makecell[{{p{\linewidth}}}]{
                \texttt{\tiny{[GM$|$GM]}}
                {'word': 'сила', 'assignment': 'team'}
            }
        }
    }
    & & \\ \\

    \theutterance \stepcounter{utterance}  
    & & \multicolumn{4}{p{0.6\linewidth}}{
        \cellcolor[rgb]{0.9,0.9,0.9}{
            \makecell[{{p{\linewidth}}}]{
                \texttt{\tiny{[GM$\rangle$P2]}}
                \texttt{Слово курча було відгадано правильно. Ваш хід закінчився.} \\
\texttt{Тепер надайте ще один список слів, розділених комами, щонайменше з 1 і щонайбільше з 1 слів із наведеного нижче списку слів, які найкраще пов'язані або найбільш тісно асоціюються зі словом 'спорт'. Пам'ятайте, що потрібно почати вашу відповідь з ВІДГАДКА: . Зауваження: деякі слова було вилучено зі списку порівняно з попередніми запитами.} \\
\\ 
\texttt{сила, машина, лепрекон, басейн, революція, лікарня, лід, вентилятор, африка, москва, екран, дошка} \\
            }
        }
    }
    & \\ \\

    \theutterance \stepcounter{utterance}  
    & & & \multicolumn{4}{p{0.6\linewidth}}{
        \cellcolor[rgb]{0.9,0.9,0.9}{
            \makecell[{{p{\linewidth}}}]{
                \texttt{\tiny{[GM$\langle$P2]}}
                \texttt{ВІДГАДКА: лід} \\
            }
        }
    }
     \\ \\

    \theutterance \stepcounter{utterance}  
    & & & \multicolumn{2}{p{0.3\linewidth}}{
        \cellcolor[rgb]{0.9,0.9,0.9}{
            \makecell[{{p{\linewidth}}}]{
                \texttt{\tiny{[GM$|$GM]}}
                ['лід']
            }
        }
    }
    & & \\ \\

    \theutterance \stepcounter{utterance}  
    & & & \multicolumn{2}{p{0.3\linewidth}}{
        \cellcolor[rgb]{0.9,0.9,0.9}{
            \makecell[{{p{\linewidth}}}]{
                \texttt{\tiny{[GM$|$GM]}}
                {'word': 'лід', 'assignment': 'innocent'}
            }
        }
    }
    & & \\ \\

    \theutterance \stepcounter{utterance}  
    & & & \multicolumn{2}{p{0.3\linewidth}}{
        \cellcolor[rgb]{0.9,0.9,0.9}{
            \makecell[{{p{\linewidth}}}]{
                \texttt{\tiny{[GM$|$GM]}}
                \texttt{лід} \\
            }
        }
    }
    & & \\ \\

    \theutterance \stepcounter{utterance}  
    & & & \multicolumn{2}{p{0.3\linewidth}}{
        \cellcolor[rgb]{0.9,0.9,0.9}{
            \makecell[{{p{\linewidth}}}]{
                \texttt{\tiny{[GM$|$GM]}}
                {'word': 'басейн', 'assignment': 'opponent'}
            }
        }
    }
    & & \\ \\

    \theutterance \stepcounter{utterance}  
    & & & \multicolumn{2}{p{0.3\linewidth}}{
        \cellcolor[rgb]{0.9,0.9,0.9}{
            \makecell[{{p{\linewidth}}}]{
                \texttt{\tiny{[GM$|$GM]}}
                {'hidden': {'team': ['революція', 'дошка', 'сила'], 'innocent': ['лепрекон', 'екран', 'лікарня', 'москва', 'вентилятор', 'машина'], 'opponent': ['африка'], 'assassin': []}, 'revealed': {'team': {'team': ['міст', 'клен', 'контракт', 'ліжко', 'відьма', 'курча'], 'innocent': ['бар', 'лід'], 'opponent': [], 'assassin': []}, 'opponent': {'team': [], 'innocent': [], 'opponent': ['лимон', 'злодій', 'стопа', 'концерт', 'прибулець', 'залізо', 'басейн'], 'assassin': []}}}
            }
        }
    }
    & & \\ \\

    \theutterance \stepcounter{utterance}  
    & & \multicolumn{4}{p{0.6\linewidth}}{
        \cellcolor[rgb]{0.9,0.9,0.9}{
            \makecell[{{p{\linewidth}}}]{
                \texttt{\tiny{[P1$\langle$GM]}}
                \texttt{Слово лід було відгадано, але це нейтральне слово. Хід вашого товариша по команді закінчився.} \\
\texttt{Тепер надайте ще одну підказку, пов'язану з деякими з ваших командних слів, що залишилися, та список пов'язаних цільових слів. Пам'ятайте: почніть підказку з 'ПІДКАЗКА: ', перейдіть на новий рядок і почніть ваш список цільових слів, розділених комами, з 'ЦІЛІ: '. Зверніть увагу: деякі слова було вилучено зі списків порівняно з попередніми запитами.} \\
\\ 
\texttt{Ваші командні слова, що залишилися: революція, дошка, сила.} \\
\texttt{Слова, що залишилися для вашого суперника: басейн, африка.} \\
\texttt{Слова{-}відволікачі, що залишилися: лепрекон, екран, лікарня, москва, вентилятор, машина.} \\
\texttt{Слова{-}вбивці, що залишилися: .} \\
            }
        }
    }
    & \\ \\

    \theutterance \stepcounter{utterance}  
    & \multicolumn{4}{p{0.6\linewidth}}{
        \cellcolor[rgb]{0.9,0.9,0.9}{
            \makecell[{{p{\linewidth}}}]{
                \texttt{\tiny{[P1$\rangle$GM]}}
                \texttt{ПІДКАЗКА: енергія} \\
\texttt{ЦІЛІ: сила} \\
            }
        }
    }
    & & \\ \\

    \theutterance \stepcounter{utterance}  
    & & & \multicolumn{2}{p{0.3\linewidth}}{
        \cellcolor[rgb]{0.9,0.9,0.9}{
            \makecell[{{p{\linewidth}}}]{
                \texttt{\tiny{[GM$|$GM]}}
                \texttt{енергія} \\
            }
        }
    }
    & & \\ \\

    \theutterance \stepcounter{utterance}  
    & & & \multicolumn{2}{p{0.3\linewidth}}{
        \cellcolor[rgb]{0.9,0.9,0.9}{
            \makecell[{{p{\linewidth}}}]{
                \texttt{\tiny{[GM$|$GM]}}
                ['сила']
            }
        }
    }
    & & \\ \\

    \theutterance \stepcounter{utterance}  
    & & & \multicolumn{2}{p{0.3\linewidth}}{
        \cellcolor[rgb]{0.9,0.9,0.9}{
            \makecell[{{p{\linewidth}}}]{
                \texttt{\tiny{[GM$|$GM]}}
                {'word': 'сила', 'assignment': 'team'}
            }
        }
    }
    & & \\ \\

    \theutterance \stepcounter{utterance}  
    & & \multicolumn{4}{p{0.6\linewidth}}{
        \cellcolor[rgb]{0.9,0.9,0.9}{
            \makecell[{{p{\linewidth}}}]{
                \texttt{\tiny{[GM$\rangle$P2]}}
                \texttt{Слово лід було відгадано, але це нейтральне слово. Ваш хід закінчився.} \\
\texttt{Тепер надайте ще один список слів, розділених комами, щонайменше з 1 і щонайбільше з 1 слів із наведеного нижче списку слів, які найкраще пов'язані або найбільш тісно асоціюються зі словом 'енергія'. Пам'ятайте, що потрібно почати вашу відповідь з ВІДГАДКА: . Зауваження: деякі слова було вилучено зі списку порівняно з попередніми запитами.} \\
\\ 
\texttt{сила, машина, лепрекон, революція, лікарня, вентилятор, африка, москва, екран, дошка} \\
            }
        }
    }
    & \\ \\

    \theutterance \stepcounter{utterance}  
    & & & \multicolumn{4}{p{0.6\linewidth}}{
        \cellcolor[rgb]{0.9,0.9,0.9}{
            \makecell[{{p{\linewidth}}}]{
                \texttt{\tiny{[GM$\langle$P2]}}
                \texttt{ВІДГАДКА: сила} \\
            }
        }
    }
     \\ \\

    \theutterance \stepcounter{utterance}  
    & & & \multicolumn{2}{p{0.3\linewidth}}{
        \cellcolor[rgb]{0.9,0.9,0.9}{
            \makecell[{{p{\linewidth}}}]{
                \texttt{\tiny{[GM$|$GM]}}
                ['сила']
            }
        }
    }
    & & \\ \\

    \theutterance \stepcounter{utterance}  
    & & & \multicolumn{2}{p{0.3\linewidth}}{
        \cellcolor[rgb]{0.9,0.9,0.9}{
            \makecell[{{p{\linewidth}}}]{
                \texttt{\tiny{[GM$|$GM]}}
                {'word': 'сила', 'assignment': 'team'}
            }
        }
    }
    & & \\ \\

    \theutterance \stepcounter{utterance}  
    & & & \multicolumn{2}{p{0.3\linewidth}}{
        \cellcolor[rgb]{0.9,0.9,0.9}{
            \makecell[{{p{\linewidth}}}]{
                \texttt{\tiny{[GM$|$GM]}}
                {'word': 'сила', 'assignment': 'team'}
            }
        }
    }
    & & \\ \\

    \theutterance \stepcounter{utterance}  
    & & & \multicolumn{2}{p{0.3\linewidth}}{
        \cellcolor[rgb]{0.9,0.9,0.9}{
            \makecell[{{p{\linewidth}}}]{
                \texttt{\tiny{[GM$|$GM]}}
                {'word': 'африка', 'assignment': 'opponent'}
            }
        }
    }
    & & \\ \\

    \theutterance \stepcounter{utterance}  
    & & & \multicolumn{2}{p{0.3\linewidth}}{
        \cellcolor[rgb]{0.9,0.9,0.9}{
            \makecell[{{p{\linewidth}}}]{
                \texttt{\tiny{[GM$|$GM]}}
                {'hidden': {'team': ['революція', 'дошка'], 'innocent': ['лепрекон', 'екран', 'лікарня', 'москва', 'вентилятор', 'машина'], 'opponent': [], 'assassin': []}, 'revealed': {'team': {'team': ['міст', 'клен', 'контракт', 'ліжко', 'відьма', 'курча', 'сила'], 'innocent': ['бар', 'лід'], 'opponent': [], 'assassin': []}, 'opponent': {'team': [], 'innocent': [], 'opponent': ['лимон', 'злодій', 'стопа', 'концерт', 'прибулець', 'залізо', 'басейн', 'африка'], 'assassin': []}}}
            }
        }
    }
    & & \\ \\

    \theutterance \stepcounter{utterance}  
    & & \multicolumn{4}{p{0.6\linewidth}}{
        \cellcolor[rgb]{0.9,0.9,0.9}{
            \makecell[{{p{\linewidth}}}]{
                \texttt{\tiny{[P1$\langle$GM]}}
                \texttt{Слово сила було відгадано правильно. Хід вашого товариша по команді закінчився.} \\
\texttt{Тепер надайте ще одну підказку, пов'язану з деякими з ваших командних слів, що залишилися, та список пов'язаних цільових слів. Пам'ятайте: почніть підказку з 'ПІДКАЗКА: ', перейдіть на новий рядок і почніть ваш список цільових слів, розділених комами, з 'ЦІЛІ: '. Зверніть увагу: деякі слова було вилучено зі списків порівняно з попередніми запитами.} \\
\\ 
\texttt{Ваші командні слова, що залишилися: революція, дошка.} \\
\texttt{Слова, що залишилися для вашого суперника: африка.} \\
\texttt{Слова{-}відволікачі, що залишилися: лепрекон, екран, лікарня, москва, вентилятор, машина.} \\
\texttt{Слова{-}вбивці, що залишилися: .} \\
            }
        }
    }
    & \\ \\

    \theutterance \stepcounter{utterance}  
    & \multicolumn{4}{p{0.6\linewidth}}{
        \cellcolor[rgb]{0.9,0.9,0.9}{
            \makecell[{{p{\linewidth}}}]{
                \texttt{\tiny{[P1$\rangle$GM]}}
                \texttt{ПІДКАЗКА: гра} \\
\texttt{ЦІЛІ: дошка} \\
            }
        }
    }
    & & \\ \\

    \theutterance \stepcounter{utterance}  
    & & & \multicolumn{2}{p{0.3\linewidth}}{
        \cellcolor[rgb]{0.9,0.9,0.9}{
            \makecell[{{p{\linewidth}}}]{
                \texttt{\tiny{[GM$|$GM]}}
                \texttt{гра} \\
            }
        }
    }
    & & \\ \\

    \theutterance \stepcounter{utterance}  
    & & & \multicolumn{2}{p{0.3\linewidth}}{
        \cellcolor[rgb]{0.9,0.9,0.9}{
            \makecell[{{p{\linewidth}}}]{
                \texttt{\tiny{[GM$|$GM]}}
                ['дошка']
            }
        }
    }
    & & \\ \\

    \theutterance \stepcounter{utterance}  
    & & & \multicolumn{2}{p{0.3\linewidth}}{
        \cellcolor[rgb]{0.9,0.9,0.9}{
            \makecell[{{p{\linewidth}}}]{
                \texttt{\tiny{[GM$|$GM]}}
                {'word': 'дошка', 'assignment': 'team'}
            }
        }
    }
    & & \\ \\

    \theutterance \stepcounter{utterance}  
    & & & \multicolumn{2}{p{0.3\linewidth}}{
        \cellcolor[rgb]{0.9,0.9,0.9}{
            \makecell[{{p{\linewidth}}}]{
                \texttt{\tiny{[GM$|$GM]}}
                \texttt{opponent has won} \\
            }
        }
    }
    & & \\ \\

\end{supertabular}
}

\end{document}
